% 第 IX 讲，第 7—8 节；原 PDF 408—426 页。
\SourcePage{408}
\SGAsection{7}{与具有半稳定约化的阿贝尔概形 $A_K$ 关联的 $S$ 上的 Raynaud 扩张 $G^\natural$；$S$ 上阿贝尔概形 $B,B'$ 的对偶性}
\textbf{7.1。}设 $S$ 是一个\Term{noetherian}{诺特}\Term{scheme}{概形}，为环 $V$ 的\Term{spectrum}{谱}，且 $V$ 对某个\Term{ideal}{理想} $\mathfrak m$ 所定义的\Term{topology}{拓扑}是\Term{separated}{分离}且\Term{complete}{完备}的。设 $A$ 是 $S$ 上的\Term{smooth}{光滑}\Term{group-scheme}{群概形}，$A^0$ 是它的\Term{neutral-component}{单位分支}；假定 $A_0/A_0^0$ 在 $S_0$ 上\Term{finite}{有限}（因而\Term{finite}{有限}\Term{etale}{平展}），并且 $A_0^0$ 是一个\Term{abelian-scheme}{阿贝尔概形}被一个\Term{torus}{环面}的\Term{extension}{扩张}。
\SourcePage{409}
仍按惯例置 $A_n=A\times_S S_n$，其中 $S_n=\operatorname{Spec}(V/\mathfrak m^{n+1})$。我们主要关心如下情形：$S$ 是一个\Term{trait}{离散赋值环谱}，\Term{generic-fiber}{泛纤维} $A_K$ 是在 $S$ 上具有\Term{semistable-reduction}{半稳定约化}的\Term{abelian-scheme}{阿贝尔概形}，而 $A$ 是 $A_K$ 的\Term{neron-model}{Néron 模型}。

设 $T_0$ 是 $A_0$ 的最大\Term{torus}{环面}，于是 $B_0=A_0^0/T_0$ 是一个\Term{abelian-scheme}{阿贝尔概形}。这时可以应用 5.1 的构造，得到\Term{formal-completion}{形式完备化} $\widehat A$ 中的一个\Term{formal-torus}{形式子环面} $\widehat T$。已知对每个 $n$，\Term{quotient}{商}
\begin{equation}\tag{7.1.1}
 B_n=A_n^0/T_n
\end{equation}
都是\Term{representable}{可表示}的（SGA 3 VI$_A$ 3.2）；更确切地说，它是 $S_n$ 上的\Term{flat}{平坦}\Term{finite-type}{有限型}\Term{scheme}{概形}，其在 $S_0$ 上的\Term{inverse-image}{逆像}为\Term{abelian-scheme}{阿贝尔概形} $B_0$。因此 $B_n$ 是 $S_n$ 上的\Term{abelian-scheme}{阿贝尔概形}；当 $n$ 变化时，各 $B_n$ 由\Term{base-change}{基变换}相互得到，从而定义一个\Term{formal-abelian-scheme}{形式阿贝尔概形}
\begin{equation}\tag{7.1.2}
 \widehat B=(B_n)_{n\geq0}.
\end{equation}
下面给出一些条件，以保证这个\Term{formal-abelian-scheme}{形式阿贝尔概形}是\Term{algebraizable}{可代数化}的（EGA III 5.4.1）；更确切地说，它是 $S$ 上某个\Term{projective}{射影}\Term{abelian-scheme}{阿贝尔概形}的\Term{formal-completion}{形式完备化}，我们将后者记为 $B$。根据所引文献，作为 $S$-\Term{scheme}{概形}，$B$ 在唯一\Term{isomorphism}{同构}意义下确定；同样由该处，$\widehat B$ 的\Term{group-structure}{群结构}唯一地来自 $B$ 上的\Term{group-structure}{群结构}。此时称\Term{abelian-scheme}{阿贝尔概形} $B$ 为\Term{scheme}{概形} $A$ 的\Term{abelian-part}{阿贝尔部分}，有时也记为 $A^{\mathrm{ab}}$。

\SourcePage{410}
\textbf{7.1.3。}根据 EGA III 5.4.5，$B$ 的\Term{projective}{射影性}意味着可以在 $\widehat B$ 上找到一个相对于 $S$ \Term{ample}{丰沛}的\Term{invertible-module}{可逆模} $\widehat M$，也就是 $B_n$ 上一列随 $n$ 变化而相互“\Term{gluing}{粘合}”的\Term{invertible-module}{可逆模} $M_n$，并且 $M_0$ \Term{ample}{丰沛}。一个自然想法是从 $G$ 上一个\Term{ample}{丰沛}\Term{invertible-module}{可逆模} $M$ 出发，证明对每个 $n$，$L$ 在 $G_n$ 上的\Term{restriction}{限制} $L_n$ 都是 $B_n$ 上某个\Term{invertible-module}{可逆模} $M_n$ 的\Term{inverse-image}{逆像}。我不知道这样的 $M_n$ 是否必然存在（而且对此深表怀疑）。我们采用稍有不同的方法，构造 $\widehat B$ 的一个\Term{polarization}{极化}，即由 $B_n$ 的彼此\Term{gluing}{粘合}的\Term{polarization}{极化} $P_n$ 构成的一个\Term{projective-system}{投射系}。这时已知 [5]，这些\Term{polarization}{极化}的两倍\Term{canonical}{典范}地定义 $B_n$ 上的（\Term{symmetric}{对称}）\Term{invertible-module}{可逆模} $M_n$；由于构造的\Term{canonical}{典范}性，它们必然彼此\Term{gluing}{粘合}，给出所求的 $\widehat M$。

我们把 $\operatorname{NS}(B_n/S_n)$ 的元素解释为 $(B_n,B_n)$ 上的\Term{symmetric}{对称}\Term{divisorial-correspondence}{除子对应} [5]，即 $(B_n,B_n)$ 被 $\mathbb G_{mS_n}$ 的\Term{biextension}{双扩张}。根据第 VIII 讲 3.5，这些\Term{biextension}{双扩张}与 $(A_n,A_n)$ 被 $\mathbb G_{mS_n}$ 的\Term{biextension}{双扩张}一一对应。为了得到这样的\Term{biextension}{双扩张}的一个相容系 $W_n$，只需从
\begin{equation}\tag{7.1.4}
 W,\quad (A,A)\text{ 被 }\mathbb G_{mS}\text{ 的对称双扩张}
\end{equation}
出发，令 $W_n$ 由 $W$ 经\Term{base-change}{基变换} $S_n\to S$ 得到。我们还须假设 $W_0$ 所定义的 $\operatorname{NS}(B_0)$ 中的元素是 $B_0$ 的一个\Term{polarization}{极化}，这样才能得到 $\widehat B$ 的\Term{polarization}{极化}。

\SourcePage{411}
\begin{SGAresult}{命题}{7.1.5}
沿用上述关于 $S,A,\widehat B$ 的记号，假定 $S$ \Term{normal}{正规}。则 $\widehat B$ \Term{algebraizable}{可代数化}，而它所来自的 $S$ 上的\Term{abelian-scheme}{阿贝尔概形} $B$ 在 $S$ 上是\Term{projective}{射影}的。
\end{SGAresult}
\begin{proof}
显然可设 $A=A^0$ 且 $S$ \Term{integral}{整}。由 M. Raynaud [23 XI 1.13]，$S$ \Term{normal}{正规}这一假设蕴涵 $A$ 在 $S$ 上\Term{quasi-projective}{拟射影}。因此取 $A$ 上一个\Term{relatively-ample}{相对丰沛}的\Term{invertible-module}{可逆模} $L$，并如第 VIII 讲 4.12 那样，在 $A\times_S A$ 上构造\Term{birigidified}{双刚化}的模
\[
 \delta(L)=m^*(L)\operatorname{pr}_1^*(L)^{-1}\operatorname{pr}_2^*(L)^{-1}.
\]
由于 $A_0^0$ 的\Term{fiber}{纤维}的\Term{unipotent-rank}{幺幂秩}为零，$A$ 的\Term{fiber}{纤维}也如此（由 3.1 容易得出），故由第 VIII 讲 7.5，$\delta(L)$ 定义 $(A,A)$ 的一个\Term{biextension}{双扩张} $W$，它显然\Term{symmetric}{对称}，而且根据第 VIII 讲 4.13（也就是本质上根据 Raynaud [23 XI 1.11]），它在 $S_0$ 的各点满足所需的\Term{ample}{丰沛性}条件。
\end{proof}

\textbf{注记 7.1.6。}不对 $S$ 作\Term{normal}{正规性}假设，也可以仿照 13.5.3 a)，化归到\Term{normal}{正规}情形，证明 $\widehat B$ 总是\Term{algebraizable}{可代数化}的。

\textbf{7.2。}现在记录以下\Term{lemma}{引理}。
\begin{SGAresult}{\Term{lemma}{引理}}{7.2.1}
设 $S$ 是一个\Term{noetherian}{诺特}环 $V$ 的\Term{spectrum}{谱}，$V$ 对某\Term{ideal}{理想} $\mathfrak m$ 所定义的\Term{topology}{拓扑}\Term{separated}{分离}且\Term{complete}{完备}。设 $C$ 为如下\Term{group-scheme}{群概形} $A$ 所成的\Term{category}{范畴}：$A$ 在 $S$ 上\Term{smooth}{光滑}\Term{finite-type}{有限型}，$A/A^0$ 是 $S$ 上的\Term{finite}{有限}\Term{etale}{平展}\Term{group-scheme}{群概形}，而 $A^0$ 是一个\Term{abelian-scheme}{阿贝尔概形} $B$ 被一个\Term{isotrivial-torus}{等平凡环面} $T$ 的\Term{extension}{扩张}（SGA 3 IX 1.1）。
\SourcePage{412}
设 $C'$ 为 $\operatorname{Specf}(A)$ 上的\Term{finite-type}{有限型}\Term{formal-scheme}{形式概形}\Term{category}{范畴}中的群所成的\Term{category}{范畴}（EGA I 10.3.3），并令 $A\mapsto\widehat A$ 为从 $C$ 到 $C'$ 的“沿 $A_0=A\times_S S_0$ 作\Term{formal-completion}{形式完备化}”\Term{functor}{函子}，其中 $S_0=\operatorname{Spec}(V/\mathfrak m)$。则：

 a) 上述\Term{functor}{函子}是\Term{fully-faithful}{全忠实}的。

 b) 设 $\mathcal A$ 是 $C'$ 的一个\Term{object}{对象}，在 $\widehat S=\operatorname{Specf}(A)$ 上\Term{smooth}{光滑}（即 $S_n=\operatorname{Spec}(A/\mathfrak m^{n+1})$ 上的普通\Term{scheme}{概形} $A_n=\mathcal A\times_{\widehat S}S_n$ 都\Term{smooth}{光滑}）。假设 $A_0/A_0^0$ 是 $S_0$ 上的\Term{finite}{有限}\Term{etale}{平展}\Term{group-scheme}{群概形}，而 $A_0^0$ 是\Term{abelian-scheme}{阿贝尔概形} $B_0$ 被\Term{isotrivial-torus}{等平凡环面} $T_0$ 的\Term{extension}{扩张}。令 $\mathcal A^0=(A_n^0)_{n\geq0}$，$\mathcal A/\mathcal A^0=(A_n/A_n^0)_{n\geq0}$。则 $\mathcal A/\mathcal A^0$ 在 $\widehat S$ 上\Term{finite}{有限}\Term{etale}{平展}，而 $\mathcal A^0$ 以本质唯一且关于 $\mathcal A$ \Term{functorial}{函子性}的方式，成为 $\widehat S$ 上的一个\Term{formal-abelian-scheme}{形式阿贝尔概形} $\mathcal B$ 被一个\Term{formal-torus}{形式环面} $\mathcal T$ 的\Term{extension}{扩张}，它们分别\Term{reduced}{约化}为 $B_0$ 和 $T_0$。$\mathcal A$ 属于上述\Term{functor}{函子} $C\to C'$ 的\Term{essential-image}{本质像}（也称其“\Term{algebraizable}{可代数化}”）的充分必要条件是 $\mathcal B$ 也如此。
\end{SGAresult}

\textbf{7.2.1.1。}先注意：在任意\Term{scheme}{概形} $S$ 上，若 $A$ 是一个\Term{smooth}{光滑}\Term{group-scheme}{群概形}，且 $A^0$ 是\Term{abelian-scheme}{阿贝尔概形}被\Term{torus}{环面}的\Term{extension}{扩张}，则这一\Term{extension}{扩张}结构本质唯一，并且以显然的意义关于 $A$ 是\Term{functorial}{函子性}的。事实上，我们的断言立即归结为：从\Term{torus}{环面} $T$ 到\Term{abelian-scheme}{阿贝尔概形} $B$ 的每个\Term{morphism}{态射}都是零。众所周知，在\Term{geometric-fiber}{几何纤维}上确实如此 [15]，再由 SGA 3 IX 5.3 即得结论。将此结果用于 7.2.1 中 $S_n$ 上的\Term{smooth}{光滑}\Term{group-scheme}{群概形} $A_n$ 所成的\Term{inductive-system}{归纳系}，
\SourcePage{413}
便得到 $\widehat S$ 上的\Term{formal-group-scheme}{形式群概形}的相应结果。此外，7.1 已经指出，若 $A_0^0$ 是\Term{abelian-scheme}{阿贝尔概形}被\Term{torus}{环面}的\Term{extension}{扩张}，则 $A_n^0$ 也如此。这就证明了 b) 的第一项断言。

\textbf{7.2.1.2。}另一方面，假设 $\widehat S$ 上的\Term{formal-group-scheme}{形式群概形} $\mathcal A$ 形如 $\widehat A$，其中 $A\in\operatorname{ob}C$。若 $A^0$ 是\Term{abelian-scheme}{阿贝尔概形} $B$ 被\Term{torus}{环面} $T$ 的\Term{extension}{扩张}，则立即看出 $\mathcal B\simeq\widehat B$，这证明 $\mathcal A$ \Term{algebraizable}{可代数化}蕴涵 $\mathcal B$ \Term{algebraizable}{可代数化}。反过来，假设 $\mathcal B$ \Term{algebraizable}{可代数化}，即来自一个\Term{abelian-scheme}{阿贝尔概形} $B$，并证明 $\mathcal A$ 也\Term{algebraizable}{可代数化}。先证明 $\mathcal A^0=(A_n^0)$ \Term{algebraizable}{可代数化}。为此注意，由第 VIII 讲 3.7，$B_n$ 被 $T_n$ 的\Term{extension}{扩张} $A_n^0$ 可以解释为\Term{homomorphism}{同态}
\begin{equation}\tag{*}
 \underline M_n\longrightarrow B_n',
\end{equation}
其中 $\underline M_n$ 是 $S_n$ 上与 $T_n$ \Term{dual}{对偶}的\Term{twisted-constant-group}{扭常值群}（SGA 3 X 5.7），$B_n'$ 是 $B_n$ 的\Term{dual-abelian-scheme}{对偶阿贝尔概形}。由于\Term{restriction}{限制}到 $S_0$ 或任一 $S_n$ 的\Term{functor}{函子}，给出 $S$ 上\Term{isotrivial-torus}{等平凡环面}的\Term{category}{范畴}与 $S_0$ 上\Term{isotrivial-torus}{等平凡环面}的\Term{category}{范畴}之间的\Term{equivalence}{等价}（SGA 3 X 3.2），各 $T_n$ 来自 $S$ 上一个在唯一\Term{isomorphism}{同构}意义下确定的\Term{torus}{环面} $T$。将其\Term{dual}{对偶}群记为 $\underline M$；其在 $S_n$ 上的\Term{restriction}{限制}\Term{canonical}{典范}\Term{isomorphism}{同构}于 $\underline M_n$。$T$ 的\Term{isotrivial}{等平凡性}可以解释为 $\underline M$ 是 $S$ 上一些\Term{finite}{有限}\Term{scheme}{概形} $\underline M^i$ 的和\Term{scheme}{概形}（SGA 3 X 5.11）。立即得到各\Term{homomorphism}{同态} (*) 来自唯一的\Term{homomorphism}{同态}
\begin{equation}\tag{**}
 \underline M\longrightarrow B',
\end{equation}
\SourcePage{414}
其中 $B'$ 是 $B$ 的\Term{dual-abelian-scheme}{对偶阿贝尔概形}。这个\Term{homomorphism}{同态}又解释为 $B$ 被 $T$ 的一个\Term{extension}{扩张} $A^0$，它显然给出 $\mathcal A^0$ 的一个\Term{algebraization}{代数化}。

\textbf{7.2.1.3。}另一方面，按假设 $A_0/A_0^0$ 是 $S$ 上的\Term{finite}{有限}\Term{group-scheme}{群概形} $\Phi_0$，故 $A_n/A_n^0$ 也都是 $S$ 上的\Term{finite}{有限}\Term{group-scheme}{群概形} $\Phi_n$（它们的\Term{representability}{可表示性}由\SGAsourcegap{原 PDF 第 414 页，第 7.2.1.3 段\Term{lie-bracket}{括号}内“leur représentabilité résultant de”之后的引注空缺。}得出），而且它们必然在 $S$ 上\Term{etale}{平展}。于是存在 $S$ 上一个在唯一\Term{isomorphism}{同构}意义下确定的\Term{etale}{平展}\Term{group-scheme}{群概形} $\Phi$，它产生各 $\Phi_n$ 所定义的\Term{formal-group-scheme}{形式群概形} $\mathcal A/\mathcal A^0$。因为 $G=A^0$ 是\Term{divisible}{可除}的，立即看出，$S$ 上的\Term{finite}{有限}\Term{locally-constant}{局部常值}群（例如 $\Phi$）被\Term{group-scheme}{群概形} $A^0$ 的 fppf \Term{sheaf}{层}\Term{extension}{扩张}所成的\Term{category}{范畴}，\Term{canonical}{典范}\Term{equivalence}{等价}于\Term{sheaf}{层}
\[
 H=\underline{\operatorname{Hom}}_{\mathrm{gr}}(\Phi,G)
\]
之下的\Term{torsor}{挠子}所成的\Term{category}{范畴}，且此\Term{equivalence}{等价}与任意\Term{base-change}{基变换}相容。将同样的描述用于 $\Phi_n$ 被 $G_n$ 的\Term{extension}{扩张}，以 $H_n\simeq\underline{\operatorname{Hom}}_{\mathrm{gr}}(\Phi_n,G_n)$ 下的\Term{torsor}{挠子}描述它们，则各 $A_n$ 定义各 $H_n$ 下\Term{torsor}{挠子}的一个相容系 $P_n$。由于 $G$ 是\Term{abelian-scheme}{阿贝尔概形}被\Term{torus}{环面}的\Term{extension}{扩张}，容易\Term{pushout}{推出} $H$ 是 $S$ 上\Term{finite}{有限}且\Term{locally-free}{局部自由}的\Term{group-scheme}{群概形}；由此容易得出，\Term{functor}{函子} $P\mapsto\widehat P=(P_n)$ 给出 $S$ 上 $H$ 下的\Term{torsor}{挠子}\Term{category}{范畴}与 $\widehat H$ 下的“\Term{form}{形式}”\Term{torsor}{挠子}\Term{category}{范畴}之间的\Term{equivalence}{等价}；后者即各 $H_n$ 下的\Term{torsor}{挠子} $P_n$ 所成的相容系。
\SourcePage{415}
因此，由\Term{extension}{扩张} $A_n$ 所定义的 $P_n$ 来自 $S$ 上同一个 $H$ 下的\Term{torsor}{挠子} $P$。由此得到 $\Phi$ 被 $G=A^0$ 的一个\Term{extension}{扩张} $A$，它“\Term{algebraization}{代数化}”了 $\mathcal A$，但还须证明它是\Term{representable}{可表示}的。可以把 $A$ 看作 $\Phi$ 上群为 $A^0$ 的一个\Term{torsor}{挠子}，并立即看出，它在\Term{finite}{有限}\Term{surjective}{满射}\Term{locally-free}{局部自由}\Term{morphism}{态射} $P\times_S\Phi\to\Phi$ 所作的\Term{base-change}{基扩张}后成为\Term{trivial}{平凡}的（因为\Term{base-change}{基变换} $P\to S$ 使所考虑的 $\Phi$ 被 $A^0$ 的\Term{extension}{扩张} $A$ \Term{split}{分裂}）。于是可用 M. Raynaud 的许多\Term{effective-descent}{有效下降}判据之一 [23 XI 3.1 6)] 得到结论。

\textbf{7.2.1.4。}这就证明了 b)，只需再证明 a) 的\Term{fully-faithful}{全忠实性}断言。留给读者验证，这已经本质上包含在 b) 的证明方法中；读者也可另寻一个直接而独立的证明。

\textbf{7.2.1.5。}应注意，若假设 $A^0$ 在 $S$ 上\Term{quasi-projective}{拟射影}，或\Term{equivalence}{等价}地，$B$ 在 $S$ 上\Term{projective}{射影}，则不必援引 Raynaud 那个精细的\Term{descent}{下降}结果。特别地，若从 $\widehat S$ 上的\Term{formal-group-scheme}{形式群概形} $\widehat A$ 出发，其中 $A$ 如 7.1.5 所述，则此条件成立；而这正是我们将使用 7.2.1 的唯一情形。

\textbf{7.2.2。}现在将 7.2.1 用于 7.1 所考虑的\Term{formal-group-scheme}{形式群概形} $\widehat A$，并取 7.1 所述的有利情形，即\Term{formal-abelian-scheme}{形式阿贝尔概形} $\widehat B$ \Term{algebraizable}{可代数化}（参见 7.1.5），此外假设 $S_0$ 上的\Term{torus}{环面} $T_0$ \Term{isotrivial}{等平凡}。于是得到 $S$ 上一个\Term{group-scheme}{群概形} $A^\natural$，称为 $A$ 在 $S$ 上的\Term{raynaud-group}{Raynaud 群}；
\SourcePage{416}
它在\Term{isomorphism}{同构}意义下，由 $\operatorname{Specf}(V)=\widehat S$ 上\Term{formal-group-scheme}{形式群概形}的一个\Term{canonical}{典范}\Term{isomorphism}{同构}
\begin{equation}\tag{7.2.3}
 (A^\natural)^\wedge\simeq\widehat A
\end{equation}
所刻画。由 7.2.1，
\begin{equation}\tag{7.2.4}
 \Phi=A^\natural/A^{\natural0}
\end{equation}
是 $S$ 上的\Term{finite}{有限}\Term{etale}{平展}群，而\Term{neutral-component}{单位分支}
\[
 A^{\natural0}=A^{0\natural}
\]
具有\Term{extension}{扩张}结构
\begin{equation}\tag{7.2.5}
 0\longrightarrow T\longrightarrow A^{\natural0}\longrightarrow B\longrightarrow0,
\end{equation}
其中 $B$ 是 $S$ 上的\Term{abelian-scheme}{阿贝尔概形}，其\Term{formal-completion}{形式完备化}为 7.1 所构造的\Term{formal-abelian-scheme}{形式阿贝尔概形} $\widehat B$，而 $T$ 则（在唯一\Term{isomorphism}{同构}意义下）刻画为 $S$ 上\Term{lifting}{提升} $T_0$ 的\Term{isotrivial-torus}{等平凡环面}。这个\Term{extension}{扩张}结构 (7.2.5) \Term{restriction}{限制}到各 $S_n$ 后，给出 (7.1.1) 所定义的 $A_n^{\natural0}=A_n^0$ 上的\Term{extension}{扩张}结构。

\textbf{7.2.6。}另一方面，由 7.2.1 立即得出，\Term{raynaud-group}{Raynaud 群} $A^\natural$ 以及\Term{exact-sequence}{正合列} (7.2.5)，都关于（满足所述条件的）$A$ \Term{functorial}{函子性}地依赖，而且它们的构造与\Term{adic}{进}\Term{topological-ring}{拓扑环}的\Term{homomorphism}{同态} $V\to V'$ 所关联的任意\Term{base-change}{基变换} $S'\to S$ 交换。

\textbf{7.3。}固定一个\Term{prime-number}{素数} $\ell$，考虑\Term{pro-ell-group}{pro-$\ell$ 群} $T_\ell(A^0)$。把 $S$ 上的\Term{finite}{有限}\Term{scheme}{概形}与相应的\Term{formal-scheme}{形式概形}等同，并如通常那样定义\Term{fixed-part}{“固定部分”} $T_\ell(A^0)^f$，
\SourcePage{417}
便得到一个如 (5.1.5) 那样的\Term{isomorphism}{同构}
\begin{equation}\tag{7.3.1}
 T_\ell(A^0)^f\simeq T_\ell(\widehat A^0).
\end{equation}
再如 (5.1.6) 那样定义\Term{toric-part}{“环面部分”}
\begin{equation}\tag{7.3.2}
 T_\ell(A^0)^t=T_\ell(\widehat T)\simeq T_\ell(T)
 \subset T_\ell(A^0)^f\simeq T_\ell(\widehat A^0),
\end{equation}
立即得到\Term{abelian-part}{“阿贝尔部分”}的\Term{isomorphism}{同构}
\begin{equation}\tag{7.3.3}
 T_\ell(A^0)^{\mathrm{ab}}\mathrel{\mathop=^{\mathrm{def}}}
 T_\ell(A^0)^f/T_\ell(A^0)^t\simeq T_\ell(\widehat B).
\end{equation}
到这里还没有用到 $\widehat B$ \Term{algebraizable}{可代数化}的假设；正是这一假设使我们能构造\Term{abelian-scheme}{阿贝尔概形} $B$ 和 $B$ 被 $T$ 的\Term{raynaud-extension}{Raynaud 扩张}。由于显然有一个\Term{canonical}{典范}\Term{isomorphism}{同构} $T_\ell(\widehat B)\simeq T_\ell(B)$（这是把 (7.3.1) 用于\Term{abelian-scheme}{阿贝尔概形} $B$ 的特例！），由 (7.3.3) 还得到
\begin{equation}\tag{7.3.4}
 T_\ell(A^0)^{\mathrm{ab}}\simeq T_\ell(B).
\end{equation}
顺便指出，(7.2.3) 与 (7.3.1) 还蕴涵\Term{canonical}{典范}\Term{isomorphism}{同构}
\begin{equation}\tag{7.3.5}
 T_\ell(A^0)^f\simeq T_\ell(A^{\natural0}),
\end{equation}
因为右端已经是一个\Term{profinite-group}{pro-有限群}，即等于它的\Term{fixed-part}{固定部分}。\Term{canonical}{典范}\Term{exact-sequence}{正合列}
\begin{equation}\tag{7.3.6}
 0\longrightarrow T_\ell(A^0)^t\longrightarrow T_\ell(A^0)^f
 \longrightarrow T_\ell(B)\longrightarrow0
\end{equation}
无论如何都能（不假设 $\widehat B$ \Term{algebraizable}{可代数化}）解释为对\Term{exact-sequence}{正合列}
\[
 0\longrightarrow\widehat T\longrightarrow\widehat A^0
 \longrightarrow\widehat B\longrightarrow0
\]
应用\Term{functor}{函子} $T_\ell$ 而得的\Term{exact-sequence}{正合列}，
\SourcePage{418}
也可以看作对普通\Term{group-scheme}{群概形}的\Term{canonical}{典范}\Term{exact-sequence}{正合列}
\[
 0\longrightarrow T\longrightarrow A^{\natural0}\longrightarrow B\longrightarrow0
\]
应用\Term{functor}{函子} $T_\ell$ 而得。

\textbf{7.4。}现在假设给定 $S$ 上第二个\Term{group-scheme}{群概形} $A'$，满足与 $G$ 相同的假设，并假设 $A,A'$ 的\Term{fiber}{纤维}\Term{connected}{连通}。设 $W$ 是 $(A,A')$ 被 $\mathbb G_{mS}$ 的\Term{biextension}{双扩张}。由此得到一个\Term{pairing}{配对}（第 VIII 讲 2.2）
\begin{equation}\tag{7.4.1}
 T_\ell(A)\times T_\ell(A')\longrightarrow
 T_\ell(\mathbb G_{mS})=\mathbb Z_\ell(1),
\end{equation}
并通过\Term{restriction}{限制}到\Term{fixed-part}{固定部分}得到\Term{pairing}{配对}
\begin{equation}\tag{7.4.2}
 T_\ell(A)^f\times T_\ell(A')^f
 =T_\ell(\widehat G)\times T_\ell(\widehat G')\longrightarrow\mathbb Z_\ell(1).
\end{equation}
它也可以解释为由 $(\widehat A,\widehat A')$ 被 $\widehat{\mathbb G}_{mS}$ 的\Term{biextension}{双扩张} $\widehat W$ 所定义的\Term{pairing}{配对}，其中 $\widehat W$ 由 $W$ 经\Term{completion}{完备化}得到（为简单起见，这里将\Term{formal-group-scheme}{形式群概形}的\Term{biextension}{双扩张}定义为 $(A_n,A_n')$ 被 $\mathbb G_{mS_n}$ 的\Term{biextension}{双扩张} $W_n$ 所成的相容族）。由第 VIII 讲 3.5，已知 $W_n$ 来自 $(B_n,B_n')$ 的一个\Term{biextension}{双扩张} $W_n^{\mathrm{ab}}$（其中 $B_n'$ 是 $G_n'$ 的\Term{abelian-part}{阿贝尔部分}），所以 (7.4.2) \Term{resolution}{分解}为
\begin{equation}\tag{7.4.3}
 T_\ell(A)^{\mathrm{ab}}\times T_\ell(A')^{\mathrm{ab}}
 =T_\ell(\widehat B)\times T_\ell(\widehat B')\longrightarrow\mathbb Z_\ell(1),
\end{equation}
它也是由 $(\widehat B,\widehat B')$ 被 $\widehat{\mathbb G}_{mS}$ 的\Term{biextension}{双扩张}
\[
 \widehat W^{\mathrm{ab}}=(W_n^{\mathrm{ab}})
\]
所定义的\Term{pairing}{配对}。当 $\widehat B,\widehat B'$ \Term{algebraizable}{可代数化}，因而有 $S$ 上的\Term{abelian-scheme}{阿贝尔概形} $B,B'$ 时，
\SourcePage{419}
由 EGA IV 5.1.4，$\widehat W^{\mathrm{ab}}$ 来自一个在唯一\Term{isomorphism}{同构}意义下确定的\Term{biextension}{双扩张}
\[
 W^{\mathrm{ab}},\quad (B,B')\text{ 被 }\mathbb G_{mS}\text{ 的双扩张},
\]
相应的\Term{pairing}{配对}
\begin{equation}\tag{7.4.4}
 T_\ell(B)\times T_\ell(B')\longrightarrow T_\ell(\mathbb G_{mS})
\end{equation}
\Term{canonical}{典范}\Term{isomorphism}{同构}于 (7.4.3)。

将 $B'$ 的\Term{dual-abelian-scheme}{对偶阿贝尔概形}记为 $D(B')$，则 $W^{\mathrm{ab}}$ 定义一个\Term{homomorphism}{同态}
\begin{equation}\tag{7.4.5}
 u:B\longrightarrow D(B'),
\end{equation}
其性质反映 (7.4.4)，亦即 (7.4.3) 的\Term{pairing}{配对}性质。例如，\Term{pairing}{配对} (7.4.3) 在一点 $t\in S$ 是\Term{perfect-duality}{完美对偶}，当且仅当 $u_t:B_t\to D(B')_t$ 是一个\Term{kernel}{核}的\Term{rank}{秩}与 $\ell$ 互素的\Term{isogeny}{同源}。这蕴涵 $u$ 本身就是\Term{isogeny}{同源}（$S$ 是\Term{local}{局部}的，因而\Term{connected}{连通}），而且 $\ker u$ 在 $S$ 上\Term{flat}{平坦}\Term{finite}{有限}，故 $\ker u$ 在每一点的\Term{rank}{秩}都相同。由此看出，刚才考虑的条件与所选的点 $t$ 无关。

\textbf{7.4.6。}特别地，这就完成了 5.4 的证明：在该\Term{theorem}{定理}的条件下（可以假定 $S$ \Term{complete}{完备}），将上述结果用于 $A^0$ 和 $A'^0$，已经证明对于每个 $\ell$，\Term{pairing}{配对} (7.4.3) 在\Term{generic-point}{泛点} $\eta$ 都是\Term{perfect-duality}{完美对偶}，所以在点 $s$ 也一样。

\SourcePage{420}
\textbf{7.4.7。}同样，由上述考虑，\Term{orthogonality-theorem}{正交性定理} 5.2 也变得清楚了，无须借助\Term{lemma}{引理} 5.2.4。事实上，在 7.4 的过程中已经证明，\Term{pairing}{配对} $(5.2.3.1)=(7.4.2)$ 在两边的\Term{toric-part}{环面部分}上都为零，从而\Term{resolution}{分解}为\Term{pairing}{配对} (7.4.3)。已知它在所考虑的\Term{pro-group}{pro-群}的\Term{special-fiber}{特殊纤维}上\Term{nondegenerate}{非退化}，只需证明它在\Term{generic-fiber}{泛纤维}上也\Term{nondegenerate}{非退化}。这正是因为，若 (7.4.5) 在\Term{special-fiber}{特殊纤维}上诱导一个\Term{isogeny}{同源}，则它就是\Term{isogeny}{同源}。

\textbf{7.4.8。}回到 7.4 的\Term{generic}{一般}条件，并假设给定一个\Term{open-subset}{开集} $U\subset S-S_0$，使得 $A_U,A_U'$ 是经\Term{biextension}{双扩张} $W_U$ 互为\Term{dual}{对偶}的\Term{abelian-scheme}{阿贝尔概形}。将\Term{exact-sequence}{正合列} (7.3.6) 及关于 $A'$ 的相应\Term{exact-sequence}{正合列}\Term{restriction}{限制}到 $U$，便在 Barsotti--Tate \Term{pro-ell-group}{pro-$\ell$ 群} $T_\ell(A_U)$ 和 $T_\ell(A_U')$ 上得到\Term{canonical}{典范}\Term{filtration}{滤过}
\begin{equation}\tag{7.4.9}
 \left\{\begin{aligned}
 T_\ell(A_U)&\supset T_\ell(A_U)^f\supset T_\ell(A_U)^t\supset0,\\
 T_\ell(A_U')&\supset T_\ell(A_U')^f\supset T_\ell(A_U')^t\supset0.
 \end{aligned}\right.
\end{equation}
我断言，这两个\Term{filtration}{滤过}对于\Term{pairing}{配对} (7.4.1) 互为\Term{dual}{对偶}（这里该\Term{pairing}{配对}是\Term{perfect-duality}{完美对偶}，见第 VIII 讲 3.2），即仍有关系式 (5.2.2) 及其\Term{symmetric}{对称}关系式成立（将其中的 $A_K,A_K'$ 换成 $A_U,A_U'$）。此外，上面已经看到，这个\Term{orthogonality}{正交性}断言也\Term{equivalence}{等价}于说，
\SourcePage{421}
$(B,B')$ 被 $\mathbb G_{mS}$ 的\Term{biextension}{双扩张} $W^{\mathrm{ab}}$ 在 $U$ 上是\Term{perfect-duality}{完美对偶}（若 $U$ 在 $S$ 中\Term{dense}{稠密}，则也\Term{equivalence}{等价}于在 $S$ 上如此）。此性质只需在 $U$ 的每一点 $\eta$ 验证；通常的论证便将问题化归到 $S$ 是以 $\eta$ 为\Term{generic-point}{泛点}的\Term{trait}{离散赋值环谱}的情形，而此情形已经知道。

另一方面，\Term{orthogonality}{正交性}关系使我们能将公式 (5.5.7)、(5.5.8)、(5.5.9) 推广到当前情形，只须将\Term{restriction}{限制}到\Term{generic-point}{泛点} $\eta$ 的运算换为\Term{restriction}{限制}到\Term{open-subset}{开集} $U$ 的运算；因此 5.6 也推广到当前情形。这里不再把上述公式和陈述在当前语境下重新写一遍。

现在从\Term{neron-model}{Néron 模型}的角度，总结关于\Term{abelian-part}{阿贝尔部分}及其\Term{pairing}{配对}所得到的结果。

\begin{SGAresult}{命题}{7.5}
设 $S$ 是一个\Term{complete}{完备}\Term{trait}{离散赋值环谱}，其\Term{function-field}{函数域}为 $K$。对 $K$ 上每个在 $S$ 上具有\Term{semistable-reduction}{半稳定约化}的\Term{abelian-scheme}{阿贝尔概形} $A_K$，将 7.1 的方法用于 $A_K$ 的\Term{neron-model}{Néron 模型} $A$ 的\Term{neutral-component}{单位分支}，便可定义它的\Term{abelian-part}{阿贝尔部分} $B=A^{\mathrm{ab}}$；它是 $S$ 上一个关于 $A_K$ \Term{functorial}{函子性}地依赖的\Term{abelian-scheme}{阿贝尔概形}。对于另一个具有\Term{semistable-reduction}{半稳定约化}的\Term{abelian-scheme}{阿贝尔概形} $A_K'$，可以给 $(A_K,A_K')$ 上每个\Term{divisorial-correspondence}{除子对应} $W_K$ 关联 $(A^{\mathrm{ab}},A'^{\mathrm{ab}})$ 上唯一的\Term{divisorial-correspondence}{除子对应} $W^{\mathrm{ab}}$，使得对每个整数 $n\geq0$，$W_n^{\mathrm{ab}}$ 沿
\[
 A_n^0\longrightarrow B_n=A_n^{0\mathrm{ab}},\qquad
 A_n'^0\longrightarrow B_n'=A_n'^{0\mathrm{ab}}
\]
的\Term{inverse-image}{逆像}为由 $W$ 经\Term{reduced}{约化}得到的 $(A_n^0,A_n'^0)$ 的\Term{biextension}{双扩张} $W_n$，其中 $W$ 是 $W_K$ 到 $(A^0,A'^0)$ 的\Term{canonical-extension}{典范延拓}（1.4）。
\SourcePage{422}
$W^{\mathrm{ab}}$ 的形成与沿\Term{abelian-scheme}{阿贝尔概形}的\Term{morphism}{态射} $\overline A_K\to A_K$、$\overline A_K'\to A_K'$ 取\Term{correspondence}{对应}的\Term{inverse-image}{逆像}相容。若 $W_K$ 是\Term{perfect-duality}{完美对偶}，则 $W^{\mathrm{ab}}$ 也如此。若 $W_K$ 是与 $A_K$ 的一个\Term{polarization}{极化}关联的 $(A_K,A_K)$ 上的\Term{divisorial-correspondence}{除子对应}，则 $W^{\mathrm{ab}}$ 也与 $A^{\mathrm{ab}}$ 的一个\Term{polarization}{极化}关联。
\end{SGAresult}
\begin{proof}
由上面对 $W^{\mathrm{ab}}$ 的刻画，其形成与\Term{inverse-image}{逆像}的相容性断言是显然的。关于\Term{polarization}{极化}的补充，在 7.1.5 证明 $\widehat B$ 的\Term{algebraizability}{可代数化性}时已经看到。
\end{proof}

\textbf{7.6。}仍假设 $S$ 是\Term{complete}{完备}\Term{trait}{离散赋值环谱}，而 $A$ 是 $K$ 上一个\Term{abelian-scheme}{阿贝尔概形} $A_K$ 的\Term{neron-model}{Néron 模型}。置
\begin{equation}\tag{7.6.1}
 A_K^\natural=(A^\natural)_K,
\end{equation}
其中 $A^\natural$ 是 Raynaud \Term{scheme}{概形}（7.2.2），即依次以\Term{abelian-scheme}{阿贝尔概形} $B$ 和\Term{torus}{环面} $T$ 对\Term{finite}{有限}\Term{etale}{平展}群 $\Phi$ 所作的\Term{extension}{扩张}。因此 $A_K^\natural$ 本身是 $\Phi_K$ 被 $A_K^{\natural0}=A_K^{0\natural}$ 的\Term{extension}{扩张}，而后者是\Term{abelian-scheme}{阿贝尔概形} $B_K$ 被\Term{torus}{环面} $T_K$ 的\Term{extension}{扩张}。对于每个\Term{prime-number}{素数} $\ell$，由 (7.3.5) 和 (7.3.2) 有
\begin{equation}\tag{7.6.2}
 T_\ell(A_K^\natural)\simeq T_\ell(A_K^{\natural0})\simeq T_\ell(A_K)^f,
 \qquad T_\ell(T_K)\simeq T_\ell(A_K)^t.
\end{equation}
\textbf{7.6.3。}由 3.3 应注意到，$A_K^{\natural0}$ 的形成与\Term{complete}{完备}\Term{trait}{离散赋值环谱}的\Term{morphism}{态射} $S'\to S$ 所作的任意\Term{base-change}{基变换}交换。

\SourcePage{423}
\SGAsection{8}{非半稳定情形的 Raynaud 扩张 $A_K^{\natural0}$ 及 ind-群 $A_K^\flat$}
\textbf{8.1。}仍假设\Term{trait}{离散赋值环谱} $S$ \Term{complete}{完备}，但不再要求 $A_K$ 具有\Term{semistable-reduction}{半稳定约化}。与通常一样，取 $K$ 的一个\Term{finite}{有限}\Term{galois-extension}{Galois 扩张} $K'$，使 $A_{K'}$ 相对于 $S$ 在 $K'$ 中的\Term{normalization}{正规化} $S'$ 具有\Term{semistable-reduction}{半稳定约化}。考虑与 $A_{K'}$ 关联的\Term{raynaud-extension}{Raynaud 扩张} $A_{K'}^\natural$（7.6.1）。显然，$K'/K$ 的\Term{galois-group}{Galois 群}通过结构转移\Term{action}{作用}于此\Term{group-scheme}{群概形}，因而由\Term{galois-descent}{Galois 下降}得到 $K$ 上一个\Term{group-scheme}{群概形} $A_{K,K'}^\natural$，它依赖于\Term{extension}{扩张} $K'$ 的选择。然而，由 7.6.3，其\Term{neutral-component}{单位分支}在\Term{canonical}{典范}\Term{isomorphism}{同构}意义下不依赖于此选择，我们将它简记为 $A_K^{\natural0}$。由\Term{descent}{下降}看出，$A_K^{\natural0}$ 是一个\Term{extension}{扩张}
\begin{equation}\tag{8.1.1}
 0\longrightarrow T_K\longrightarrow A_K^{\natural0}\longrightarrow B_K\longrightarrow0,
\end{equation}
其中 $T_K$ 是\Term{torus}{环面}，$B_K$ 是 $K$ 上的\Term{abelian-scheme}{阿贝尔概形}。同样，由 $T_\ell(A_K)$ 的\Term{essentially-fixed-part}{本质固定部分}与\Term{essentially-toric-part}{本质环面部分}的定义（4.1 和 6.5），\Term{isomorphism}{同构} (7.6.2) 给出\Term{canonical}{典范}\Term{isomorphism}{同构}
\begin{equation}\tag{8.1.2}
 T_\ell(A_K^\natural)\simeq T_\ell(A_K)^{ef},\qquad
 T_\ell(T_K)\simeq T_\ell(A_K)^{et},
\end{equation}
取\Term{quotient}{商}便有
\begin{equation}\tag{8.1.3}
 T_\ell(B_K)\simeq T_\ell(A_K)^{ef}/T_\ell(A_K)^{et}.
\end{equation}

\SourcePage{424}
这样，\Term{filtration}{滤过}
\begin{equation}\tag{8.1.4}
 T_\ell(A_K)\supset T_\ell(A_K)^{ef}\supset T_\ell(A_K)^{et}\supset0
\end{equation}
中的两个\Term{quotient}{商}就用\Term{group-scheme}{群概形}的\Term{extension}{扩张} (8.1.1) 确定了。为了确定第三个\Term{quotient}{商} $T_\ell(A_K)/T_\ell(A_K)^{ef}$，注意由 (6.5.2) 有一个\Term{canonical}{典范}\Term{isomorphism}{同构}
\begin{equation}\tag{8.1.5}
 T_\ell(A_K)/T_\ell(A_K)^{ef}\simeq D\bigl(T_\ell(A_K')^{et}\bigr),
\end{equation}
其中 $D=\underline{\operatorname{Hom}}(-,T_\ell(\mathbb G_m))$ 表示 Cartier \Term{dual}{对偶}（在 Barsotti--Tate \Term{pro-group}{pro-群}的意义下），$A_K'$ 表示 $A_K$ 的\Term{dual-abelian-scheme}{对偶阿贝尔概形}。于是考虑与 $K$ 上的\Term{abelian-scheme}{阿贝尔概形} $A_K'$ 关联的、形如 (8.1.1) 的\Term{exact-sequence}{正合列}
\begin{equation}\tag{8.1.6}
 0\longrightarrow T_K'\longrightarrow A_K'^{\natural0}\longrightarrow B_K'\longrightarrow0
\end{equation}
以及相应的形如 (8.1.2) 的\Term{isomorphism}{同构}，公式 (8.1.5) 就给出
\begin{equation}\tag{8.1.7}
 T_\ell(A_K)/T_\ell(A_K)^{ef}\simeq D\bigl(T_\ell(T_K')\bigr)
 \simeq\underline M_K\otimes_{\mathbb Z}\mathbb Z_\ell,
\end{equation}
其中置（推广 (5.5.2) 所引入的记号）
\begin{equation}\tag{8.1.8}
 \underline M_K=D(T_K'),\qquad \underline M_K'=D(T_K),
\end{equation}
而这里 $D$ 表示普通\Term{group-scheme}{群概形}意义下的 Cartier \Term{dual}{对偶} $\underline{\operatorname{Hom}}(-,\mathbb G_{mK})$。因此 $\underline M_K$ 和 $\underline M_K'$ 是 $K$ 上的\Term{twisted-constant-group}{扭常值群}，在 $K'$ 上成为\Term{constant-group}{常值群}，并\Term{isomorphism}{同构}于形如 $\mathbb Z^r$ 的群（$r=\dim T=\dim T'$）。

\SourcePage{425}
\textbf{8.2。}条件与 8.1 相同。显然，\Term{extension}{扩张} (8.1.1) 关于\Term{abelian-scheme}{阿贝尔概形} $A_K$ \Term{functorial}{函子性}地依赖，其形成与任意\Term{trait}{离散赋值环谱}\Term{morphism}{态射} $\widetilde S\to S$ 交换。选取 $A_K$ 的一个\Term{polarization}{极化}，从而得到\Term{isogeny}{同源} $A_K\to A_K'$，由此得到从 (8.1.1) 到 (8.1.6) 的一个\Term{homomorphism}{同态}；（考察 $T_\ell$）立即看出，它逐项都是\Term{isogeny}{同源}。同样，利用 7.5 和\Term{descent}{下降}可见，$K$ 上的\Term{abelian-scheme}{阿贝尔概形} $A_K,\widetilde A_K$ 被 $\mathbb G_{mK}$ 的每个\Term{biextension}{双扩张}，都定义 $(B_K,\widetilde B_K)$ 被 $\mathbb G_{mK}$ 的一个\Term{biextension}{双扩张}，其中 $B_K,\widetilde B_K$ 分别表示与 $A_K,\widetilde A_K$ 对应的\Term{raynaud-extension}{Raynaud 扩张}的\Term{abelian-quotient}{阿贝尔商}；当 $W$ 是\Term{perfect-duality}{完美对偶}时，$(B_K,\widetilde B_K)$ 的相应\Term{biextension}{双扩张}也如此。特别地，沿用 7.6 的记号，\Term{abelian-scheme}{阿贝尔概形} $B_K,B_K'$ \Term{canonical}{典范}地互为\Term{dual}{对偶}。当 $A_K,\widetilde A_K$ 变化时，$(A_K,\widetilde A_K)$ 上的\Term{biextension}{双扩张}与 $(B_K,\widetilde B_K)$ 上的\Term{biextension}{双扩张}之间的对应具有显然的\Term{functorial}{函子性}；这些性质可由 7.5 所述的相应性质经\Term{descent}{下降}得到，其陈述留给读者。上述映射与任意\Term{trait}{离散赋值环谱}\Term{base-change}{基变换} $\widetilde S\to S$ 的相容性，也留给读者。

\textbf{8.3。}回到群 $A_{K,K'}^\natural$，研究它关于 $K'$ 的变换性。显然，$K$ 的两个\Term{isomorphism}{同构}\Term{extension}{扩张} $K',K_1'$ 给出\Term{canonical}{典范}\Term{isomorphism}{同构}的群，因此只须考虑 $K$ 的\Term{separable-closure}{可分闭包} $\overline K$ 中的子\Term{extension}{扩张} $K'$。若 $K'\subset K''$ 是两个这样的 $K$ 的\Term{finite}{有限}\Term{galois-extension}{Galois 子扩张}，则由 3.3 立即得出，
\SourcePage{426}
有一个\Term{canonical}{典范}\Term{open-immersion}{开浸入}
\begin{equation}\tag{8.3.1}
 A_{K,K'}^\natural\lhook\joinrel\longrightarrow A_{K,K''}^\natural,
\end{equation}
由此可将第一个群等同于第二个群的一个\Term{open-subgroup}{开子群}。当 $K'$ 变化时，各 $A_{K,K'}^\natural$ 构成 $K$ 上\Term{smooth}{光滑}\Term{finite-type}{有限型}群的一个\Term{inductive-system}{归纳系}，其转移\Term{morphism}{态射}是\Term{open-immersion}{开浸入}。若置
\begin{equation}\tag{8.3.2}
 \Phi_{K,K'}=A_{K,K'}^\natural/A_{K,K'}^{\natural0},
\end{equation}
则各 $\Phi_{K,K'}$ 本身构成 $K$ 上\Term{finite}{有限}\Term{etale}{平展}群的一个\Term{inductive-system}{归纳系}，后面将更详细地研究它。

我们置
\begin{equation}\tag{8.3.3}
 A_K^\flat=\varinjlim_{K'}A_{K,K'}^\natural,
\end{equation}
此\Term{limit}{极限}在 $K$ 上 fppf \Term{sheaf}{层}的\Term{category}{范畴}中取。由于所考虑\Term{inductive-system}{归纳系}中的转移\Term{morphism}{态射}都是\Term{open-immersion}{开浸入}，显然此\Term{limit}{极限}实际上由 $K$ 上一个\Term{group-scheme}{群概形}表示（\Term{generic}{一般}不\Term{quasi-compact}{拟紧}），我们仍然用同一个记号 $A_K^\flat$ 表示它。注意，对于这个\Term{group-scheme}{群概形}有
\begin{equation}\tag{8.3.4}
 A_K^{\flat0}=A_K^{\natural0},\qquad
 A_K^\flat/A_K^{\flat0}\simeq\varinjlim_{K'}\Phi_{K,K'}.
\end{equation}
最后这个群将在 11.10 中确定：它\Term{canonical}{典范}\Term{isomorphism}{同构}于 $\underline M_K\otimes\mathbb Q/\mathbb Z$（采用 (8.1.8) 的记号）。
